\honest{Empty Labels}{Eliezer Yudkowsky}{2008}

Consider, yet again, Aristotelian categories. An object has properties A, B, C, D and E. I use letters, I tell Fred, because in Aristotelian logic it makes no difference what the properties are or what I call them.

I invent ``zawa'' for objects that are A, C and D; ``yokie'' for B and E; ``xippo'' for E but not D. I tell Fred, ``Object 1 is zawa and yokie, but not xippo,'' and when I spell it out, Fred asks, ``Amazing! You can tell all that just by looking?'' Then come ``bolo'' (A, C and yokie), ``mun'' (A, C and xippo) and ``merlacdonian'' (bolo and mun). Pointlessly confusing, I agree. Replace each label by its definition and ``merlacdonian'' becomes two nested lists of letters. The list is the ``entire information'' of the label; the rules work without any label at all.

Define ``human'' as a mortal featherless biped, and the classic syllogism reads: all mortal featherless bipeds are mortal; Socrates is a mortal featherless biped; therefore Socrates is mortal. The label hid the premise and made the conclusion look new; replacing it with the definition exposes a tautology that tells you nothing about the world. You cannot call Socrates human until you have seen him die. \nb{The objection is Mill's, from 1843; see ``The Parable of Hemlock.''}

There is an idea, which you may have noticed I hate, that ``you can define a word any way you like.'' It came from the Aristotelian notion of categories, because an agent following those rules exactly would never reach a contradiction. It would not reach much of anything. \nb{No source is given for this history, and the idea is older than Aristotle. In Plato's \textsc{Cratylus}, Hermogenes says that any name you give ``is the right one, and if you change that and give another, the new name is as correct as the old.''}

The labels are not so much arbitrary as unnecessary: ``The labels are only there to create the illusion of inference.'' \nb{Earlier the post had called a label ``a human convenience (or inconvenience).''} So the proverb should not be ``I can define a word any way I like,'' nor ``Defining a word never has any consequences,'' but ``Definitions don't need words.''
